% sec_09_interbody
\section{Inter-body reflection}\label{sec:interbody}

The master surface law of \cref{sec:surface-law} captures first-bounce
absorption only. At each illuminated point a fraction $1-T_0$ of the
incident power is reflected rather than absorbed, and on a non-convex
body part of that reflected power can re-illuminate another body part:
across the gap between the legs, between an arm and the torso, or under
the chin. This section formalises the single specular recapture as an
optional second-order correction. It is off by default. The production
law omits it because the validated body-averaged enhancement stays
inside the diffraction and tissue-property error bars.

\subsection{Reflected power and the specular direction}

At a front-facing point $\rr$ with incidence cosine
$\mu=\nhat\cdot(-\khat)$ the Fresnel reflection coefficients are the
siblings of the transmission coefficients of \cref{sec:fresnel},
\begin{equation}\label{eq:r-coeffs}
  r_s(\mu)=\frac{\mu-\zeta}{\mu+\zeta},\qquad
  r_p(\mu)=\frac{\ntilde^2\mu-\zeta}{\ntilde^2\mu+\zeta},\qquad
  \zeta=\sqrt{\ntilde^2-1+\mu^2},\ \ \RE\zeta>0,
\end{equation}
and they satisfy the polarisation-wise energy identity
$T_s+\abs{r_s}^2=1$, $T_p+\abs{r_p}^2=1$ at the lossy interface.

\begin{definition}[Reflected power fraction]\label{def:rho-refl}
With TE and TM power weights $w_s,w_p$ ($w_s+w_p=1$) carried from the
incident polarisation as in \cref{sec:pol}, the locally reflected
fraction of the incident power density is
\begin{equation}\label{eq:rho-refl}
  \rho_R(\rr)=\abs{r_s(\mu)}^2 w_s+\abs{r_p(\mu)}^2 w_p
            =1-T_0(\rr),
\end{equation}
the Fresnel reflectance. The specularly reflected ray leaves along
\begin{equation}\label{eq:k-refl}
  \khat_{\mathrm{refl}}=\khat-2(\khat\cdot\nhat)\,\nhat
                       =\khat+2\mu\,\nhat,
\end{equation}
the mirror of $\khat$ about the tangent plane.
\end{definition}

\subsection{One-bounce recapture}

The single specular recapture casts the reflected ray
\cref{eq:k-refl} from every lit triangle and reuses the visibility
acceleration structure of \cref{sec:geom} to test whether it
strikes another body triangle before escaping. When it does, the
reflected power density $\rho_R\,\Sinc$ is deposited at the hit point,
weighted by its own incidence cosine, and added to the direct
first-bounce term there. This is the $\mathtt{inter\_body=specular1}$
axis: a single optional pass, orthogonal to the diffraction and
polarisation flags, that never feeds back into the visibility bake.

To bound the cumulative effect of all bounces, let $f(\rr)$ denote the
fraction of power reflected from $\rr$ that strikes the body again.
Summing the geometric series of successive bounces enhances the
absorbed power at a point by
\begin{equation}\label{eq:radiosity-C}
  C(f)=\frac{1}{1-\Tbar_R\,f},\qquad \Tbar_R=1-\Tbar\approx 0.46,
\end{equation}
with $\Tbar_R$ the flux-weighted reflectance. Under diffuse
(Lambertian) reflection the recaptured fraction is bounded by the
cosine-weighted hemisphere blocked by other body parts,
\begin{equation}\label{eq:f-bound}
  f(\rr)\le 1-\Omega(\rr),
\end{equation}
where $\Omega$ is exactly the ambient-occlusion exposure fraction of
\cref{sec:geom}. For a convex body every reflected ray escapes,
$f\equiv 0$, and the first-bounce law is exact.

\begin{proposition}[Self-compensation]\label{prop:self-comp}
The enhancement $C$ is largest precisely where $\Omega$ is smallest. A
deep concavity sees little direct power but recaptures a large share of
the reflected power, so the product $\Omega\,C$ stays close to $\Omega$.
By \cref{eq:f-bound} a point with $\Omega=0.3$ admits $f\le 0.7$, hence
$C\le 1/(1-0.46\times 0.7)\approx 1.47$ and $\Omega\,C\le 0.44$,
recovering less than a fifth of the $70\%$ shadow deficit.
\end{proposition}

\subsection{Energy bookkeeping}

The recapture is an exchange of power between body parts inside the
incidence-operator balance derived for the coherent regime in
\cref{sec:coherent}. Writing $\Qop$ for the absorption operator and
$\Qop_{\mathrm{ref}}$ for its reflection complement (the Gram of the
specularly scattered surface field), the landed power partitions as
\begin{equation}\label{eq:op-balance}
  \Qop_{\mathrm{inn}}=\Qop+\Qop_{\mathrm{ref}},
\end{equation}
where $\Qop_{\mathrm{inn}}$ maps a precoder to the power that lands on
the body irrespective of whether it is then absorbed or scattered
away. The single specular recapture promotes a slice of
$\Qop_{\mathrm{ref}}$ back into $\Qop$ to second order in the body's
own albedo $\bar\alpha=\operatorname{tr}\Qop_{\mathrm{ref}}/
\operatorname{tr}\Qop_{\mathrm{inn}}$, leaving evanescent surface
waves, contour creeping, and triple-and-higher bounces in the
residual.

\subsection{Validated bounds}

Measured on the Thelonious phantom (\num{23826} triangles) over a sweep
of incidence directions, the correction is small but the monograph's
original specular estimate was optimistic.

\begin{remark}[What the data say]\label{rem:interbody-data}
Three findings hold across every concavity region.
(i) The body-averaged enhancement is $C\approx 1.05$ to $1.08$, at most
$8\%$, well below the diffraction (${\sim}10\%$) and tissue-property
(${\sim}10\%$) uncertainties, which justifies omitting reflection from
the default law.
(ii) The specular recapture is \emph{comparable} to the diffuse bound,
not one third of it: in a concavity the mirror direction points back
at the opposing surface, giving
$f_{\mathrm{spec}}\approx 1.08\,f_{\mathrm{diffuse}}$, so the specular
enhancement matches the diffuse $C$ rather than the near-unity value
once assumed.
(iii) The deepest concavity reaches a local $C\approx 1.85$. This
matters for spatial maps in tight gaps, not for compliance, which is
set by the most-illuminated surface where $f\to 0$.
\end{remark}

The correction is therefore retained as the optional
$\mathtt{specular1}$ pass for spatial-map fidelity and recorded in the
error budget of \cref{sec:validation}, but it is excluded from the
default equations.
